\chapter{Fermions in one dimension and bosonization}\label{chap:bosonization}

The theories of interest contain fermions, but the library simulates bosons. In one spatial dimension this is
not an approximation: a massless Dirac fermion \emph{is} a compact massless boson, a statement known as
bosonization~\cite{Coleman1975,Mandelstam1975,Giamarchi2003}. This chapter explains the part of the dictionary
the library uses and checks it on the spectrum of free fermions on a ring.

\section{Chiral fermions}
In one dimension a massless fermion field splits into a left-moving and a right-moving component,
$\psi = (\psi_L, \psi_R)$, with
\begin{equation}
  H = \int dx\,\Bigl[\psi_R^\dagger(-i\partial_x)\psi_R + \psi_L^\dagger(i\partial_x)\psi_L\Bigr].
\end{equation}
A right mover has energy $E = k$ for momentum $k$, a left mover $E = -k$; both move at the speed of light and
never scatter into each other without interactions. Each component has its own particle number,
$N_L = \int\psi_L^\dagger\psi_L$ and $N_R = \int\psi_R^\dagger\psi_R$ (measured relative to the filled Dirac sea).
A \emph{chiral} theory treats them differently, for example by giving them different gauge charges.

\paragraph{Spectrum on a ring.} On a ring of length $L$ (with antiperiodic boundary conditions for the fermions)
the single-particle momenta are $(r+\frac12)\,2\pi/L$. Adding $N_R$ right movers above the Fermi sea fills the
lowest levels and costs
\begin{equation}
  \frac{2\pi}{L}\sum_{r=0}^{N_R-1}\Bigl(r+\tfrac12\Bigr) = \frac{2\pi}{L}\,\frac{N_R^2}{2},
\end{equation}
the same for holes ($N_R<0$) and for left movers. The lowest energy at given fermion numbers is therefore
\begin{equation}\label{eq:fermion-sector}
  E(N_L, N_R) = \frac{2\pi}{L}\,(h + \bar h),\qquad h = \frac{N_L^2}{2},\quad \bar h = \frac{N_R^2}{2},
\end{equation}
and its momentum is $P = \frac{2\pi}{L}(\bar h - h)$ (left movers carry negative momentum). Above these states
there are particle-hole excitations with energies in multiples of $2\pi/L$. The numbers $h,\bar h$ are the
\emph{conformal weights} of the state; $\Delta = h + \bar h$ is its scaling dimension and $s = h - \bar h$ its
spin. A single fermion has $\Delta = |s| = \frac12$.

\section{The bosonization dictionary}
Introduce a scalar field $\varphi$ with momentum $\Pi$ such that the vector (total) and axial (left minus right)
fermion densities are
\begin{equation}\label{eq:dictionary}
  \rho_V = \psi_L^\dagger\psi_L + \psi_R^\dagger\psi_R = \frac{1}{\sqrt\pi}\,\partial_x\varphi,\qquad
  \rho_A = \psi_L^\dagger\psi_L - \psi_R^\dagger\psi_R = \frac{1}{\sqrt\pi}\,\Pi .
\end{equation}
The Hamiltonian of the free fermion becomes the free massless boson,
$H = \int\frac12[\Pi^2 + (\partial_x\varphi)^2]$, and the particle-hole excitations of the fermions are the
oscillator modes of the boson. It is convenient to define the \emph{dual field}
\begin{equation}
  \vartheta(x) = \int^x dy\,\Pi(y), \qquad \rho_A = \frac{1}{\sqrt\pi}\partial_x\vartheta ,
\end{equation}
so that $\varphi$ counts vector charge and $\vartheta$ axial charge in the same way.

\paragraph{Fermion numbers are zero modes.} Integrating \eqref{eq:dictionary} over the ring,
\begin{equation}
  n \equiv N_L + N_R = \frac{1}{\sqrt\pi}\bigl[\varphi(L) - \varphi(0)\bigr],\qquad
  m \equiv N_L - N_R = \frac{1}{\sqrt\pi}\int_0^L \Pi .
\end{equation}
The vector fermion number is the \emph{winding} of $\varphi$ and the axial number its \emph{total momentum}.
Since fermion numbers are integers, $\varphi$ must be compact with period $\sqrt\pi$ (\cref{sec:ring}), and $n$
and $m$ are integers with
\begin{equation}\label{eq:parity}
  n \equiv m \pmod 2,
\end{equation}
because $N_L = (n+m)/2$ and $N_R = (n-m)/2$ are integers. This parity rule is what makes the compact boson
fermionic: $(n,m) = (1,0)$ would be half a left and half a right fermion.

\paragraph{Check.} The zero-mode energy of the boson \eqref{eq:zero-mode-energy} with $W = \sqrt\pi n$,
$P_{\rm tot} = \sqrt\pi m$ is
\begin{equation}
  E = \frac{\pi(n^2+m^2)}{2L} = \frac{\pi}{L}\bigl(N_L^2 + N_R^2\bigr) = \frac{2\pi}{L}\Bigl(\frac{N_L^2}{2}
  + \frac{N_R^2}{2}\Bigr),
\end{equation}
exactly the fermionic result \eqref{eq:fermion-sector}. The lattice momentum $P = -\frac1a\sum_j P_j\ell_j$
gives
\begin{equation}
  s = -\frac{PL}{2\pi} = \frac{nm}{2} = \frac{N_L^2 - N_R^2}{2} = h - \bar h .
\end{equation}
The library reproduces both exactly on any ring
(\code{analysis.sectors.sector_energy}; tested for several $(N_L,N_R)$). Note the sign convention: left movers have
negative momentum and positive spin, $s = h - \bar h$.

\section{Many fermions and their charges}
A theory contains left movers $L_1,\dots,L_{n_L}$ and right movers $R_1,\dots,R_{n_R}$, with charges
$Q_{L,i}$, $Q_{R,j}$ under a $U(1)$ gauge field (for $U(1)^k$: one row of charges per gauge field). To bosonize,
the library pairs $L_i$ with $R_i$ into a Dirac fermion and describes each pair by one compact boson
$\varphi_i$. The gauge current $\sum_i (Q_{L,i}\rho_{L,i} + Q_{R,i}\rho_{R,i})$ becomes
\begin{equation}\label{eq:va-charges}
  \sum_i \bigl(Q^V_i \rho_{V,i} + Q^A_i \rho_{A,i}\bigr),\qquad
  Q^V = \frac{Q_L + Q_R}{2},\quad Q^A = \frac{Q_L - Q_R}{2}.
\end{equation}
The pairing is bookkeeping: physical results must not depend on it, and the tests check this explicitly for the
two possible pairings of the 3-4-5-0 model. If $n_L \neq n_R$, the missing partners are added as neutral fermions
(charge $0$), which are free and decoupled (\code{ChiralTheory.pairing}).

\begin{python}
from latsimpkg import ChiralTheory

t = ChiralTheory(QL=[3, 4], QR=[5, 0], g=0.3)
p = t.pairing()
p.QV, p.QA          # [[4, 2]], [[-1, 2]]
t.pairing(order=[1, 0]).QV   # the other pairing: [[1.5, 4.5]]
\end{python}
